\documentclass[10pt,a4paper]{amsart}
\usepackage[margin=19mm]{geometry}
\usepackage{amsmath,amssymb,mathtools}
\usepackage[scr=rsfs,cal=euler]{mathalfa}
\usepackage{xfrac}
\usepackage[unicode,hidelinks,bookmarks=false]{hyperref}
\newcommand{\ie}{i.e.\ }
\newcommand{\cf}{cf.\ }
\newcommand{\resp}{respectively\ }
\newcommand{\Subsection}{\subsection}
\providecommand{\nobreakdash}{\nobreak\hbox{-}}
\setlength{\emergencystretch}{2em}
\multlinegap0pt
\usepackage{bm}
\usepackage{bbm}
\usepackage{dsfont}
\usepackage{xspace}
\usepackage{cancel}
\usepackage{mathtools}
\usepackage{amsthm}
\makeatletter
\@ifundefined{thm}{\newtheorem{thm}{Thm}}{}
\@ifundefined{prop}{\newtheorem{prop}[thm]{Proposition}}{}
\makeatother
\newcommand{\sk}{\smallskip}
\title[A story of webs: the webs by conics on del Pezzo quartic sur]{A story of webs: the webs by conics on del Pezzo quartic surfaces and Gelfand-MacPherson's web of the spinor tenfold}
\author{Luc Pirio}
\date{Source: arXiv:2507.12180v1 (CC BY 4.0)}
\begin{document}
\maketitle

\noindent\emph{Source and licence.} A story of webs: the webs by conics on del Pezzo quartic surfaces and Gelfand-MacPherson's web of the spinor tenfold. Version arXiv:2507.12180v1 (CC BY 4.0), distributed under the Creative Commons Attribution 4.0 International licence, as recorded in the arXiv licence metadata for this exact version and shown on the abstract page https://arxiv.org/abs/2507.12180v1. Full rights notice: licenses/arxiv-2507.12180-CC-BY-4.0.txt.\par\smallskip

\noindent\textit{Editorial scope.} Counted unit: the Prop whose source label is \texttt{Prop:Transformation-sigma-i} in the article (source0.tex lines 3384--3402, source label \texttt{Prop:Transformation-sigma-i}), delivered together with the article's own complete original proof of it (source0.tex lines 3404--3463, one \texttt{proof} environment of 60 lines). No mathematics is written, repaired or re-derived: statement and proof are the article's own text line for line. Required context retained uncounted: none. Every definition, display and label of those lines is retained unchanged. Omitted from this package and not redistributed: 1--3383, 3403--3403, 3464--8933. Nothing inside a retained excerpt is omitted or altered: every retained body is delivered exactly as the article's lines are written, so no omission receipt of any kind accompanies this package. Citations: the retained excerpts carry no citation command at all, so no bibliography entry of the article is transcribed and no reference is redistributed. Reference adaptation: none was needed and none was made; every reference inside the delivered unit resolves inside it, so no unresolved reference or citation is produced. Printed slips kept literal: no printed word, symbol, hypothesis or number of the counted statement or of its proof was altered. Layout and class: the article's own class is \texttt{amsart} with packages including dynkin-diagrams, tikz, psfrag, listings, bigints, multirow, oubraces, fontenc, lmodern, xcolor, cellspace, slashbox, color, xy; the wrapper is the lane's pinned \texttt{amsart} editorial wrapper at 10pt on A4, which restates the theorem-like environments the retained text needs and adds the article's own macro definitions that the retained text needs (\texttt{\textbackslash sk}), with extra wrapper packages (bm, bbm, dsfont, xspace, cancel, mathtools). The delivered numbering is the unit's own. Assets: no retained span contains a figure environment, a graphics-inclusion command, a PDF-page inclusion, a table or a TikZ picture; no image file is copied into this package, the package declares no asset dependency and no image file is redistributed with it. Licence: version arXiv:2507.12180v1 of this article is distributed under the Creative Commons Attribution 4.0 International licence, as recorded in the arXiv licence metadata for this exact version and shown on the abstract page https://arxiv.org/abs/2507.12180v1. A full rights notice is delivered at licenses/arxiv-2507.12180-CC-BY-4.0.txt. Characters: every retained excerpt is pure ASCII, so the delivered engine, XeLaTeX with its default Latin Modern text font, reports no missing character.
\begin{prop}
\label{Prop:Transformation-sigma-i}
%%%%%%%%%%%%%%%%%%%%%%%%%%%%%%%%%%%%%%
{\bf 1.} As elements of $\boldsymbol{AR}^2\big( \boldsymbol{W}^{GM}_{\boldsymbol{Y}_5} \big)$, one has $\sigma_i^*
\big(  {\bf HLOG}^\omega_{ \boldsymbol{Y}_5}\big)=-  \, {\bf HLOG}^\omega_{ \boldsymbol{Y}_5}$
for any $i=1,\ldots,5$.  Hence 
$ {\bf HLOG}^\omega_{ \boldsymbol{Y}_5}$ spans a 1-dimensional non-trivial $W_{D_5}$-subrepresentation 
of $\boldsymbol{AR}^2\big( \boldsymbol{W}^{GM}_{\boldsymbol{Y}_5} \big)$ 
which necessarily is the signature representation. 
\sk 

%%%%%%%%%%%%%%%%%%%%%%%%%%%%%%%%%%%%%%
\noindent 
{\bf 2.} There is a similar statement for the linear action of $W_{D_5}$ on 
${\rm Gr}^\bullet  \boldsymbol{\mathcal A\mathcal R}^2$: the complex line 
${\rm Gr}^1  \boldsymbol{\mathcal A\mathcal R}^2$ spanned by ${\bf HLOG}_{ \boldsymbol{Y}_5}$ is $W_{D_5}$-stable and is the signature representation. 
%{\bf 3.}  The abelian relation $ {\bf HLOG}_{ {\mathcal Y}_5}$ is stable 
%%%%%%%%%%%%%%%%%%%%%%%%%%%%%%%%%%%%%%
\end{prop}

\begin{proof}%%%%%%%%%%%%%%%%%%%%%%%%%%%%%%%
%%%%%%%%%%%%%%%%%%%%%%%%%%%%%%%%%%%%%%
For $i=1,\ldots,5$, one 
sets $ {\bf HLOG}_{i,+}^{\omega}=
\big({U_{i,1}}\big)^*\big( \Omega^\omega
\big)$  and $
{\bf HLOG}_{i,-}^{\omega}=
\big({U_{i,2}}\big)^*\big( \Omega^\omega
\big)
$. 
Then by direct elementary computations, using the explicit expressions of the $\sigma_i$'s given above, 
one easily establish the following transformation formulas: 
%%%%%%%%%%%%%%%%%%%%%%%%%%%%%%%%%%%%%%
\begin{itemize}
%%%%%%%%%%%%%%%%%%%%%%%%%%%%%%%%%%%%%%
\item for $k\in \{1,2,3,4\}$, one sets $\nu_k=(k,k+1)\in \mathfrak S_5$. Then 
for $ i=1,\ldots,5$  and 
$\epsilon=\pm$, 
one has: 
 %%%%%%%%%%%%%%%%%%%%%%%%%%%%%%%%%%%%%%
\begin{equation}
\label{Eq:Transfoformulas_1}
 %%%%%%%%%%%%%%%%%%%%%%%%%%%%%%%%%%%%%%
\sigma_k^*\big( {\bf HLOG}_{i,\epsilon}^{\omega} \big) =- 
{\bf HLOG}_{\nu_k(i),\epsilon}^{\omega} \, ; 
%\qquad \mbox{ for } i=1,\ldots,5 \mbox{ and } \, 
%\epsilon=\pm\, .
 %%%%%%%%%%%%%%%%%%%%%%%%%%%%%%%%%%%%%%
\end{equation}
 %%%%%%%%%%%%%%%%%%%%%%%%%%%%%%%%%%%%%%
\item let $\nu_5$ be the transposition exchanging $4$ and $5$ (i.e.\,$\nu_5=\nu_4$). Then 
%for $i\in \big[\hspace{-0.1cm}\big[ 10
%\big]\hspace{-0.1cm}\big]$, one sets $i^*=i+5$ if $i \leq 5$ and $i^*=i-5$ otherwise.  
%Then 
for $j\in \{1,2,3\}$, $\ell\in \{4,5\}$ and $\epsilon=\pm$, 
one has:
 %%%%%%%%%%%%%%%%%%%%%%%%%%%%%%%%%%%%%%
\begin{equation}
\label{Eq:Transfoformulas_2}
 %%%%%%%%%%%%%%%%%%%%%%%%%%%%%%%%%%%%%%
\sigma_5^*\big( {\bf HLOG}_{j,\epsilon}^{\omega} \big)=-
{\bf HLOG}_{j,\epsilon}^{\omega} 
\qquad \mbox{ and } \qquad 
\sigma_5^*\big( {\bf HLOG}_{\ell,\epsilon}^{\omega} \big)=-
{\bf HLOG}_{\nu_5(\ell),-\epsilon}^{\omega} \, 
 %%%%%%%%%%%%%%%%%%%%%%%%%%%%%%%%%%%%%%
\end{equation}
 %%%%%%%%%%%%%%%%%%%%%%%%%%%%%%%%%%%%%%
%%%%%%%%%%%%%%%%%%%%%%%%%%%%%%%%%%%%%%
\end{itemize}
%%%%%%%%%%%%%%%%%%%%%%%%%%%%%%%%%%%%%%
It follows that $\langle {\bf HLOG}_{\boldsymbol{Y}_5}^{\omega}\rangle$ is $W_{D_5}$-stable. Since this 1-dimensional representation is not trivial, it has to be the signature. 

The proof of the second point of the proposition 
%{\bf 2} 
is essentially similar. 
%The second point of the proposition follows immediately from the formulas 
%\eqref{Eq:Transfoformulas_1} and \eqref{Eq:Transfoformulas_2}.
%%%%%%%%%%%%%%%%%%%%%%%%%%%%%%%%%%%%%%
\end{proof}%%%%%%%%%%%%%%%%%%%%%%%%%%%%%%%

\end{document}
